
\section{\BookRecordIndex{T0380}隐函数定理}\label{sec:A-3}

设 $X,Y$ 是巴拿赫空间，$f:X\to Y$ 是光滑映射。对每个 $x\in X$，
以 $df(x):X\to Y$ 表示 $f$ 在 $x$ 处的微分。若此算子是双射，
则由开映射定理，其逆 $df(x)^{-1}:Y\to X$ 是有界线性算子。
逆函数定理断言：在每个使 $df(x)$ 可逆的点 $x$ 附近，$f$ 都有局部逆。
用 $B_r(x;X)$ 表示巴拿赫空间 $X$ 中以 $x$ 为中心、半径为 $r$ 的开球。
当空间由上下文确定时，简记 $B_r(x):=B_r(x;X)$、$B_r:=B_r(0;X)$。

\begin{theorem}[逆函数定理]\label{thm:A-3-1}
设 $X,Y$ 是巴拿赫空间，$U\subset X$ 是开集，$f:U\to Y$ 连续可微。
设 $x_0\in U$，且微分 $df(x_0):X\to Y$ 是双射。
则存在 $x_0$ 的开邻域 $U_0\subset U$，使 $f$ 在 $U_0$ 上的限制为单射，
$V_0:=f(U_0)$ 是 $Y$ 中的开集，$f^{-1}:V_0\to U_0$ 连续可微，且对
$y\in V_0$ 有
\[
  df^{-1}(y)=\bigl(df(f^{-1}(y))\bigr)^{-1}.
\]
因此，若 $f$ 为 $C^\ell$ 类（$\ell$ 为正整数），则 $f^{-1}$ 也是 $C^\ell$ 类。
\end{theorem}

证明基于以下关于导数接近恒等映射的 $\psi:X\to X$ 的引理。

\begin{lemma}\label{lem:A-3-2}
设 $\gamma<1$ 与 $R$ 是正实数，$X$ 是巴拿赫空间，$x_0\in X$，
$\psi:B_R(x_0)\to X$ 连续可微，且对每个 $x\in B_R(x_0)$ 有
\[
  \|\operatorname{id}-d\psi(x)\|\le\gamma.
\]
则 $\psi$ 是单射，并将 $B_R(x_0)$ 映到 $X$ 中的开集，满足
\sourcefootnote{原书下式右端的球心印作 $x_0$，但在平移映射
$\psi(x)=x+a$（其导数为恒等映射）下即不成立；证明中的归一化步骤及左端球心
均表明右端应为 $\psi(x_0)$，故据此补正。}
\[
  B_{R(1-\gamma)}(\psi(x_0))
  \subset\psi(B_R(x_0))
  \subset B_{R(1+\gamma)}(\psi(x_0)).
\]
此外，$\psi^{-1}:\psi(B_R(x_0))\to B_R(x_0)$ 连续可微，且
\begin{equation}\label{eq:A-3-1}
  d\psi^{-1}(y)=d\psi(\psi^{-1}(y))^{-1}.
\end{equation}
因此，若 $\psi$ 为 $C^\ell$ 类（$\ell$ 为正整数），则 $\psi^{-1}$
也是 $C^\ell$ 类。
\end{lemma}

\begin{proof}
用映射 $x\mapsto\psi(x_0+x)-\psi(x_0)$ 替换 $\psi$，
可假设 $x_0=0$ 且 $\psi(x_0)=0$。证明分四步。

\noindent\textbf{第 1 步：$\psi$ 是到其像上的同胚，且
$\psi(B_R)\subset B_{R(1+\gamma)}$。}
考虑映射
\[
  \phi:=\operatorname{id}-\psi:B_R\to X.
\]
对每个 $x\in B_R$，它满足 $\|d\phi(x)\|\le\gamma$。因此，
对所有 $x_1,x_2\in B_R$ 有
\begin{equation}\label{eq:A-3-2}
  \|\phi(x_1)-\phi(x_2)\|\le\gamma\|x_1-x_2\|.
\end{equation}
于是由三角不等式，对所有 $x_1,x_2\in B_R$ 有
\begin{equation}\label{eq:A-3-3}
  (1-\gamma)\|x_1-x_2\|
  \le\|\psi(x_1)-\psi(x_2)\|
  \le(1+\gamma)\|x_1-x_2\|.
\end{equation}
右侧不等式表明 $\psi(B_R)\subset B_{R(1+\gamma)}$；
左侧不等式表明 $\psi$ 是单射，且 $\psi^{-1}$ 是（Lipschitz）连续的。

\noindent\textbf{第 2 步：$B_{R(1-\gamma)}\subset\psi(B_R)$。}
取 $y\in B_{R(1-\gamma)}$，由
$\|y\|=:(1-\gamma)(R-\varepsilon)$ 定义 $\varepsilon>0$。
在式~\eqref{eq:A-3-2} 中取 $x_2=0$，可知对 $x\in B_R$ 有
$\|\phi(x)\|\le\gamma\|x\|$，从而
\[
  \|x\|\le R-\varepsilon
  \quad\Longrightarrow\quad
  \|\phi(x)+y\|\le R-\varepsilon.
\]
因此，映射 $x\mapsto\phi(x)+y$ 是半径为 $R-\varepsilon$ 的闭球上的压缩映射。
由压缩映射原理，它有唯一不动点 $x$，满足 $\|x\|\le R-\varepsilon$。
此不动点满足 $\psi(x)=x-\phi(x)=y$，故 $y\in\psi(B_R)$。

\noindent\textbf{第 3 步：$\psi(B_R)$ 是开集。}
设 $y=\psi(x)$，其中 $x\in B_R$，并选取 $\varepsilon>0$，使
$B_\varepsilon(x)\subset B_R$。由第 2 步，
\[
  B_{\varepsilon(1-\gamma)}(\psi(x))
  \subset\psi(B_\varepsilon(x))\subset\psi(B_R).
\]

\noindent\textbf{第 4 步：$\psi^{-1}$ 连续可微，其导数由式~\eqref{eq:A-3-1} 给出。}
设 $y_0=\psi(x_0)$，其中 $x_0\in B_R$，记 $D:=d\psi(x_0)$。
于是 $\|\operatorname{id}-D\|\le\gamma$，故 $D$ 可逆，并且
\[
  D^{-1}=\sum_{k=0}^{\infty}(\operatorname{id}-D)^k,
  \qquad \|D^{-1}\|\le\frac1{1-\gamma}.
\]
我们证明 $\psi^{-1}$ 在 $y_0$ 处可微，且
$d\psi^{-1}(y_0)=D^{-1}$。
\sourcefootnote{原书此处写“$\psi$ 在 $y_0$ 处可微”，但下一式及结论均讨论
$\psi^{-1}$ 的微分，且 $y_0$ 属于逆映射的定义域；据此补正主语。}
为此，固定任意 $\varepsilon>0$。由于 $\psi$ 在 $x_0$ 处可微，存在
$\delta>0$，使得对每个 $x\in B_R$，
\[
  \|x-x_0\|<\frac\delta{1-\gamma}
  \quad\Longrightarrow\quad
  \|\psi(x)-\psi(x_0)-D(x-x_0)\|
  \le\varepsilon(1-\gamma)^2\|x-x_0\|.
\]
必要时缩小 $\delta$，由第 3 步可假设
$B_\delta(y_0)\subset\psi(B_R)$。现在设
$\|y-y_0\|<\delta$，记 $x:=\psi^{-1}(y)\in B_R$。
由式~\eqref{eq:A-3-3} 得
\[
  \|x-x_0\|\le\frac1{1-\gamma}\|y-y_0\|
  \le\frac\delta{1-\gamma}.
\]
因此
\[
\begin{aligned}
  \|\psi^{-1}(y)-\psi^{-1}(y_0)-D^{-1}(y-y_0)\|
  &=\|D^{-1}(y-y_0-D(x-x_0))\|\\
  &\le\frac1{1-\gamma}
   \|\psi(x)-\psi(x_0)-D(x-x_0)\|\\
  &\le\varepsilon(1-\gamma)\|x-x_0\|\\
  &\le\varepsilon\|y-y_0\|.
\end{aligned}
\]
所以 $\psi^{-1}$ 在 $y_0$ 处可微，其导数由式~\eqref{eq:A-3-1} 给出。
由第 1 步与式~\eqref{eq:A-3-1}，映射
$y\mapsto d\psi^{-1}(y)$ 关于 $\mathcal L(X)$ 上的范数拓扑连续。
这证明了引理~\ref{lem:A-3-2}。
\end{proof}

\begin{proof}[定理~\ref{thm:A-3-1} 的证明]
不失一般性，设 $x_0=0$ 且 $f(0)=0$。考虑映射
$\psi:B_\delta(0;X)\to X$，
\[
  \psi(x):=D^{-1}f(x),
\]
其中 $D:=df(x_0)$。取 $c>0$ 使 $c\ge\|D^{-1}\|$。
\sourcefootnote{原书在此证明中直接使用常数 $c$，未先定义；
由紧接着的算子范数估计及末尾的球半径，可取任意 $c>0$ 且
$c\ge\|D^{-1}\|$。}
其微分满足
\[
  \operatorname{id}-d\psi(x)
  =\operatorname{id}-D^{-1}df(x)
  =D^{-1}(D-df(x)),
\]
因此，缩小 $\delta>0$ 后，对每个 $x\in B_\delta(0;X)$ 有
\[
  \|\operatorname{id}-d\psi(x)\|
  \le c\|D-df(x)\|\le\frac12.
\]
在引理~\ref{lem:A-3-2} 中取 $R=\delta$、$\gamma=1/2$，
可知 $\psi$ 在 $B_\delta(0;X)$ 上有连续可微的逆，
且 $\psi(B_\delta(0;X))$ 是包含 $B_{\delta/2}(0;X)$ 的开集。
所以
\[
  f(B_\delta(0;X))
  =D\psi(B_\delta(0;X))
  \supset DB_{\delta/2}(0;X)
  \supset B_{\delta/(2c)}(0;Y),
\]
且 $f$ 的逆由
\[
  f^{-1}(y)=\psi^{-1}(D^{-1}y)
\]
给出。它连续可微，公式
$df^{-1}(y)=df(f^{-1}(y))^{-1}$ 由链式法则得到。
定理~\ref{thm:A-3-1} 得证。
\end{proof}

巴拿赫空间之间的光滑映射 $f:X\to Y$ 称为\emph{弗雷德霍姆映射}，
如果对每个 $x\in X$，其微分 $df(x):X\to Y$ 都是\BookRecordIndex{T0128}弗雷德霍姆算子。
由于\BookRecordIndex{T0129}弗雷德霍姆指标在小扰动下不变，$df(x)$ 的指标与 $x$ 的选取无关，
记为 $\operatorname{index}(f)$。对任何光滑映射 $f:X\to Y$，
无论它是否为弗雷德霍姆映射，向量 $y\in Y$ 称为 $f$ 的\emph{\BookRecordIndex{T0409}正则值}，
如果对每个 $x\in f^{-1}(y)$，算子 $df(x):X\to Y$ 都是满射且有\BookRecordIndex{T0392}右逆。
\BookRecordIndex{T0380}隐函数定理断言：$f$ 的每个\BookRecordIndex{T0409}正则值 $y$ 的原像 $f^{-1}(y)$ 都是光滑流形。
此外，若 $f$ 是弗雷德霍姆映射，则 $f^{-1}(y)$ 有限维，
其维数等于 $f$ 的\BookRecordIndex{T0129}弗雷德霍姆指标。

\begin{theorem}[\BookRecordIndex{T0380}隐函数定理]\label{thm:A-3-3}
设 $X,Y$ 是巴拿赫空间，$U\subset X$ 是开集，$\ell$ 是正整数。
若 $f:U\to Y$ 为 $C^\ell$ 类且 $y$ 是 $f$ 的\BookRecordIndex{T0409}正则值，则
\[
  \mathcal M:=f^{-1}(y)\subset X
\]
是 $C^\ell$ \BookRecordIndex{T0053}巴拿赫流形，并且对每个 $x\in\mathcal M$，有
\[
  T_x\mathcal M=\ker df(x).
\]
因此，若 $f$ 是弗雷德霍姆映射，则 $\mathcal M$ 有限维，且
\[
  \dim\mathcal M=\operatorname{index}(f).
\]
\end{theorem}

证明之前，先把这个几何陈述改写为更便于分析的形式。不失一般性，设 $y=0$。
若 $x_0\in\mathcal M$，则由假设，算子
\[
  D:=df(x_0):X\to Y
\]
为满射，并有\BookRecordIndex{T0392}右逆 $Q:Y\to X$，满足 $DQ=\operatorname{id}_Y$。
\BookRecordIndex{T0392}右逆存在，当且仅当存在分解
\[
  X=\ker D\oplus\operatorname{im}Q.
\]
这里 $\ker D$ 由线性化方程 $df(x_0)\xi=0$ 的解组成；
我们预期，完整非线性方程 $f(x)=0$ 的解空间在 $x_0$ 附近局部看起来像
$\ker D$。\BookRecordIndex{T0380}隐函数定理把这一点精确化。它断言 $\ker D$ 是
$\mathcal M=f^{-1}(0)$ 在 $x_0$ 处的切空间，等价于存在 $C^\ell$ 类映射
\sourcefootnote{原书此处称所构造的 $\phi$ 为“smooth”，但本定理只假设
$f$ 为 $C^\ell$ 类，后面的证明也只得出 $\phi$ 为 $C^\ell$ 类；
故按本定理的\BookRecordIndex{T0408}正则性补正。}
\[
  \phi:\ker D\to Y,
\]
满足 $d\phi(0)=0$，且当 $x$ 充分接近 $x_0$ 时，$f(x)=0$
当且仅当存在充分小的向量 $\xi$，使得
\[
  x=x_0+\xi+Q\phi(\xi),\qquad D\xi=0.
\]
因此，记 $x_1:=x_0+\xi$，我们要寻找某个
$\eta\in\operatorname{im}Q$，使 $f(x_1+\eta)=0$。
目前只知道 $x=x_1$ 是方程 $f(x)=0$ 的近似解；
我们并不假设 $f(x_0)=0$。
\sourcefootnote{原书此前明确取 $x_0\in\mathcal M=f^{-1}(0)$，
故在这段局部几何讨论中已有 $f(x_0)=0$。此句保留原文；
它只能理解为转入下一个命题的更一般设置，那里 $x_0$ 不必为零点。}
不过，下一命题说，只要 $f(x_1)$ 充分小，陪集
$C=x_1+\operatorname{im}Q$ 中就有真正的解。

\begin{proposition}\label{prop:A-3-4}
设 $X,Y$ 是巴拿赫空间，$U\subset X$ 是开集，$f:U\to Y$ 连续可微。
设 $x_0\in U$，使 $D:=df(x_0):X\to Y$ 为满射，且有一个（有界线性）
\BookRecordIndex{T0392}右逆 $Q:Y\to X$。选取正常数 $\delta,c$，使得
$\|Q\|\le c$、$B_\delta(x_0;X)\subset U$，且
\begin{equation}\label{eq:A-3-4}
  \|x-x_0\|<\delta
  \quad\Longrightarrow\quad
  \|df(x)-D\|\le\frac1{2c}.
\end{equation}
假设 $x_1\in X$ 满足
\begin{equation}\label{eq:A-3-5}
  \|f(x_1)\|<\frac\delta{4c},
  \qquad \|x_1-x_0\|<\frac\delta8.
\end{equation}
则存在唯一的 $x\in X$，满足
\[
  f(x)=0,\qquad x-x_1\in\operatorname{im}Q,
  \qquad \|x-x_0\|<\delta.
\]
此外，$\|x-x_1\|\le2c\|f(x_1)\|$。
\end{proposition}

\begin{proof}
定义映射 $h:U\to Y$ 为
\[
  h(x):=f(x)-D(x-x_1).
\]
于是
\begin{equation}\label{eq:A-3-6}
  x+Qh(x)=x_1
  \quad\Longleftrightarrow\quad
  f(x)=0,\quad x-x_1\in\operatorname{im}Q.
\end{equation}
为说明这一点，先设
$x+Qf(x)-QD(x-x_1)=x_1$，则 $x-x_1\in\operatorname{im}Q$。
因而 $QD(x-x_1)=x-x_1$，所以 $Qf(x)=0$。
由于 $Q$ 是单射，得 $f(x)=0$；反向同理。

因此，问题归结为在 $B_\delta(x_0;X)$ 中求解 $x+Qh(x)=x_1$。
把该方程写成 $\psi(x)=x_1$，其中 $\psi:U\to X$ 定义为
\[
  \psi(x):=x+Qh(x)=x+Q(f(x)-D(x-x_1)).
\]
我们将证明 $\psi$ 满足引理~\ref{lem:A-3-2} 的假设，且
$x_1\in\psi(B_\delta(x_0;X))$。

具体地，$d\psi(x)-\operatorname{id}=Q(df(x)-D)$，
故由式~\eqref{eq:A-3-4}，
\[
  \|x-x_0\|<\delta
  \quad\Longrightarrow\quad
  \|\operatorname{id}-d\psi(x)\|
  \le c\|df(x)-D\|\le\frac12.
\]
由引理~\ref{lem:A-3-2}，$\psi$ 把 $B_\delta(x_0;X)$
双射到 $X$ 中的一个开集，并且
\sourcefootnote{原书此处右端印作
$\psi(B_{2\delta}(x_0;X))$，但 $\psi$ 未必在整个
$B_{2\delta}(x_0;X)$ 上有定义；由引理~\ref{lem:A-3-2}
的上界应为以 $\psi(x_0)$ 为中心、半径 $3\delta/2$ 的球。}
\[
  B_{\delta/2}(\psi(x_0);X)
  \subset\psi(B_\delta(x_0;X))
  \subset B_{3\delta/2}(\psi(x_0);X).
\]
再由式~\eqref{eq:A-3-3} 和~\eqref{eq:A-3-5}，
\[
\begin{aligned}
  \|x_1-\psi(x_0)\|
  &=\|\psi(x_1)-\psi(x_0)-Qf(x_1)\|\\
  &\le2\|x_1-x_0\|+c\|f(x_1)\|<\delta/2.
\end{aligned}
\]
故存在唯一的 $x\in B_\delta(x_0;X)$，使 $\psi(x)=x_1$，
等价地，$x+Qh(x)=x_1$。由式~\eqref{eq:A-3-6}，
我们找到 $x\in x_1+\operatorname{im}Q$，使 $f(x)=0$。
另外，引理~\ref{lem:A-3-2} 证明中的式~\eqref{eq:A-3-3} 给出
\[
  \|x-x_1\|
  \le2\|\psi(x)-\psi(x_1)\|
  =2\|Qf(x_1)\|
  \le2c\|f(x_1)\|.
\]
这证明了命题~\ref{prop:A-3-4}。
\end{proof}

\begin{remark}\label{rem:A-3-5}
命题~\ref{prop:A-3-4} 的证明可以用 Newton--Picard 迭代来解释。
给定使 $f(x_1)$ 很小的 $x_1$，我们要寻找
$g(\eta):=f(x_1+Q\eta)=0$ 在 $\eta=0$ 附近的解 $\eta\in Y$。
直接的 Newton--Picard 迭代方案是递归定义
\[
  \eta_{i+1}:=\eta_i-dg(0)^{-1}g(\eta_i),
  \qquad \eta_1:=0.
\]
令 $x_i:=x_1+Q\eta_i$。由于
$dg(0)=DQ=\operatorname{id}_Y$\sourcefootnote{原书在此把
$dg(0)$ 写成 $DQ=\operatorname{id}_Y$，但按
$g(\eta)=f(x_1+Q\eta)$ 直接求导，应为 $dg(0)=df(x_1)Q$；
$D=df(x_0)$ 仅在 $x_1=x_0$ 或另有条件时才与之相同。
后文迭代式可视为用冻结导数 $D$ 的迭代；这里保留原论证，不把它冒充精确的 Newton 步。}，
这个公式便是
\[
  x_{i+1}:=x_i-Qf(x_i).
\]
沿命题~\ref{prop:A-3-4} 的证明追溯到引理~\ref{lem:A-3-2}
证明中第 2 步所用的压缩映射原理，论证归结为寻找映射
$x\mapsto x_1-Qh(x)$ 的不动点，其中仍取
$h(x):=f(x)-D(x-x_1)$。这个不动点问题的迭代为
\[
  x_{i+1}=x_1-Qh(x_i)
  =x_1+QD(x_i-x_1)-Qf(x_i).
\]
又因 $x_i-x_1\in\operatorname{im}Q$，有
$x_i=x_1+QD(x_i-x_1)$。因此，这正是上面所述的迭代。
\end{remark}

\begin{proof}[定理~\ref{thm:A-3-3} 的证明]
不失一般性，设 $y=0$。取 $x_0\in U$ 使 $f(x_0)=0$，并记
$D:=df(x_0)$。由假设，$D$ 有\BookRecordIndex{T0392}右逆 $Q:Y\to X$。
按命题~\ref{prop:A-3-4} 选取常数 $c,\delta$；必要时缩小 $\delta$，
还可假设
\[
  \|\xi\|<\delta/8
  \quad\Longrightarrow\quad
  \|f(x_0+\xi)-f(x_0)-df(x_0)\xi\|<2\|\xi\|/c.
\]
取 $\xi\in\ker D$ 且 $\|\xi\|<\delta/8$。则
$x_1:=x_0+\xi$ 满足式~\eqref{eq:A-3-5}。
由命题~\ref{prop:A-3-4}，存在唯一元素 $x=x(\xi)\in X$，满足
\[
  f(x)=0,\quad x-x_0-\xi\in\operatorname{im}Q,
  \quad \|x-x_0\|<\delta.
\]
由于 $Q$ 是单射，存在唯一的 $\phi(\xi)\in Y$，使得
\[
  x=x_0+\xi+Q\phi(\xi).
\]
命题~\ref{prop:A-3-4} 还给出
$\|Q\phi(\xi)\|\le2c\|f(x_0+\xi)\|<\delta/2$。

下面证明 $\phi$ 是 $C^\ell$ 类。定义 $\psi:B_\delta(x_0)\to X$ 为
\[
  \psi(x):=Qf(x)+x-x_0-QD(x-x_0).
\]
于是
\[
  \psi(x_0)=0,
  \qquad d\psi(x)-\operatorname{id}=Q(df(x)-D).
\]
由式~\eqref{eq:A-3-4}，每个 $x\in B_\delta(x_0)$ 都满足
$\|d\psi(x)-\operatorname{id}\|\le1/2$。
因此，由引理~\ref{lem:A-3-2}，$\psi$ 在 $B_{\delta/2}(0)$
上有 $C^\ell$ 类的逆。若 $\xi\in\ker D$ 且 $\|\xi\|<\delta/8$，
则 $\|\xi+Q\phi(\xi)\|<5\delta/8<\delta$，且
$\psi(x_0+\xi+Q\phi(\xi))=\xi$，所以
\[
  x_0+\xi+Q\phi(\xi)=\psi^{-1}(\xi).
\]
因 $DQ=\operatorname{id}_Y$，这意味着
\[
  \phi(\xi)=D\psi^{-1}(\xi)-Dx_0,
  \qquad \xi\in\ker D,\quad\|\xi\|<\delta/8.
\]
故 $\phi$ 如所断言为 $C^\ell$ 类；又因
$d\psi(x_0)=\operatorname{id}$，有 $d\phi(0)=0$。

此外，若 $x\in X$ 满足 $f(x)=0$，且
$\|x-x_0\|<\delta/[8(1+c\|D\|)]<\delta$，
则可写成
\[
  x=x_0+\xi+Q\eta,\quad
  \eta:=D(x-x_0),\quad
  \xi:=x-x_0-Q\eta\in\ker D.
\]
由 $\|\eta\|\le\|D\|\|x-x_0\|$，可得
\[
  \|\xi\|\le(1+c\|D\|)\|x-x_0\|<\delta/8.
\]
因此 $\eta=\phi(\xi)$，从而 $x=\psi^{-1}(\xi)$。
于是映射
\[
  \psi:f^{-1}(0)\cap
  B_{\delta/[8(1+c\|D\|)]}(x_0)\longrightarrow\ker D
\]
是 $f^{-1}(0)$ 上的一张坐标图。转移映射显然具有 $C^\ell$ \BookRecordIndex{T0408}正则性。
\sourcefootnote{原书此处称转移映射“smooth”，但本定理只假设 $f$ 为
$C^\ell$ 类，且刚构造的坐标图及逆映射只保证 $C^\ell$；
故这里按结论所需的\BookRecordIndex{T0408}正则性表述。}
这证明了定理~\ref{thm:A-3-3}。
\end{proof}

在巴拿赫空间中，$D$ 为满射并不保证它有\BookRecordIndex{T0392}右逆。
\BookRecordIndex{T0392}右逆存在，当且仅当 $\ker D$ 在 $X$ 中有补空间
（见下面的命题~\ref{prop:A-4-1}）。特别地，每个满射\BookRecordIndex{T0128}弗雷德霍姆算子都有\BookRecordIndex{T0392}右逆。
这一点还可推广到由
\[
  (D\oplus L)(x,z)=Dx+Lz
\]
定义的算子 $D\oplus L:X\oplus Z\to Y$，其中 $D:X\to Y$ 为\BookRecordIndex{T0128}弗雷德霍姆算子。

\begin{lemma}\label{lem:A-3-6}
设 $D:X\to Y$ 是\BookRecordIndex{T0128}弗雷德霍姆算子，$L:Z\to Y$ 是有界线性算子，
且 $D\oplus L:X\oplus Z\to Y$ 是满射。则 $D\oplus L$ 有\BookRecordIndex{T0392}右逆。
此外，投影 $\Pi:\ker(D\oplus L)\to Z$ 是\BookRecordIndex{T0128}弗雷德霍姆算子，
且 $\ker\Pi\cong\ker D$、
$\operatorname{coker}\Pi\cong\operatorname{coker}D$，从而
\[
  \operatorname{index}\Pi=\operatorname{index}D.
\]
\end{lemma}

\begin{proof}
选取 $\ker D$ 在 $X$ 中的一个补空间 $X_1$，再选取有限多个向量
$z_\nu\in Z$（$\nu=1,\ldots,N$），使 $Lz_\nu$ 张成
$\operatorname{im}D$ 在 $Y$ 中的一个补空间。
则 $D\oplus L$ 的一个\BookRecordIndex{T0392}右逆是算子
\[
  Y\longrightarrow X\oplus Z:
  y\longmapsto\left(x,\sum_{\nu=1}^N\lambda_\nu z_\nu\right),
\]
其中 $x$ 与 $\lambda_1,\ldots,\lambda_N$ 选得满足
\[
  x\in X_1,\qquad
  y=Dx+\sum_{\nu=1}^N\lambda_\nu Lz_\nu.
\]
现在 $\ker\Pi=\ker D\oplus\{0\}$，且
$\operatorname{im}\Pi=L^{-1}(\operatorname{im}D)$。
因此
\[
  \frac{Z}{\operatorname{im}\Pi}
  =\frac{Z}{L^{-1}(\operatorname{im}D)}
  \cong\frac{\operatorname{im}L}
  {\operatorname{im}D\cap\operatorname{im}L}
  \cong\frac{Y}{\operatorname{im}D}.
\]
第二个同构由 $L$ 诱导；最后一个同构之所以存在，是因为
$\operatorname{im}D+\operatorname{im}L=Y$。
\end{proof}

\begin{exercise}\label{ex:A-3-7}
设 $X$ 是无限维希尔伯特空间。与有限维情形不同，
在单位球外为零的 $C^\infty$ 光滑函数
$\phi:X\to\mathbb R$ 不一定有界（尽管 $X$ 的每个点都有一个使函数有界的邻域）。
构造一个支集位于单位球内的无界光滑函数。
提示：单位球内有无穷多个两两不交、半径均为 $1/4$ 的球。
\end{exercise}
